\pagebreak

\chapter{Building 3PL}

The 3PL compiler is predominantly a Java program therefore the logical choice
of build system seemed to be Apache Ant which was designed from the ground up
for building Java programs.

The original "make" based build system had trouble dealing with dependencies
for java inner-classes and handling the contents of Java jar files. It was
hoped that Apache Ant might handle these things better (however, it is
unclear whether this was achieved).

A "make" Makefile is still supported but all it does is call Apache Ant with
any make targets passed as Ant targets.

The 3PL compiler will build on any of the three major operating systems in
use currently, namely Linux, MacOS or Windows.

The available Ant targets (and therefore "make" targets) are:
\begin{description}[leftmargin=!,labelwidth=\widthof{\bfseries versionfilter}]
\item[api]           run javadoc to create html version of API docs in doc/api
\item[build]         runs jar target then creates new hierarchy in build dir
\item[clean]         removes most output except jar file
\item[compile]       creates properties files and compiles all .java files
\item[install]       runs build target into a system directory
\item[jar]           runs compile target then builds the .jar file
\item[javacc]        rebuilds parser source from .jj file - into build/jcc-out
\item[main]          main target - runs clean target then run target
\item[release]       runs clean and install targets then creates a release
\item[revision]      print 3pl revision from git
\item[run]           runs jar target then executes jar file in place
\item[source]        creates a source archive from the current release branch
\item[version]       print 3pl version with modified flag
\item[versionfilter] replace version tokens in stdin
\item[versiontag]    print 3pl version as git tag
\end{description}
To build the compiler, simply run the command \bverb:ant build:. This will
create a subdirectory \\ \bverb:build/<version>:, where \bverb:<version>: is the
current 3PL compiler version (\bverb:<version>: will be suffixed with a
capital \bverb:M: if any files have been modified in the local git repository).

The build subdirectory will be populated with all the normal hierarchy that
would be expected if installed in the system (i.e. \bverb:bin:, \bverb:include:
and \bverb:lib:).

To install the hierarchy in a more centralised system location, run the
\bverb:ant install: command, which will install the hierarchy (including
the \bverb:<version>: subdirectory) into \bverb:/opt/3pl: (the location can
be overridden by defining the ant property \bverb:3pl.instdir: on the command
line - e.g. \\ \bverb:ant -D3pl.instdir=/path install:).

To run the compiler "in place", i.e. without building or installing the
full hierarchy, use the command \bverb:ant run:. This builds only the
required \bverb:jar: file, then runs it inside the \bverb:ant: program
(which is itself a Java program) supplying the \bverb:sources/include:
subdirectory as the include file search path.

The \bverb:ant: program is very verbose in its output. To make it run more
quietly, use the command line arguments \bverb:-q -S:. The \bverb:Makefile:
does this by default.
